\chapter{Rekonstruktion}
Nach XX entspricht die im Fernfeld aufgenommene Welle der Fouriertransformierten der Austrittswelle. Könnte sowohl Amplitude als auch Phase aufgenommen werden, so ließe sich durch eine einfache Rücktransformation die Austrittswelle rekonstruieren und somit Informationen über das Objekt gewinnen. Jedoch wird mittels eines Sensors nur die Intensität der Welle aufgezeichnet, die Informationen über die Phase gehen somit verloren. Das Problem, aus dieser unvollständigen Messung die Austrittswelle zu rekonstruieren wird als das Phasenproblem bezeichnet. Zur Lösung existiert u.A. der Ansatz der Freiflug-Holographie sowie die Methode der iterativen Phasenrekonstruktion
\section{Holographie}
Angenommen neben dem Objekt, über das Informationen gewonnen werden sollen befindet sich ein weiteres (im Weiteren als Referenz bezeichnetes) Objekt im Strahl, sodass die eingehende Welle an beiden gestreut wird. Sind Objekt und Referenz Näherungsweise 2D-Objekte in einer gemeinsamen Ebene, so gilt für die Austrittswelle $\phi_1$
\begin{equation}
	\phi_1=O+R
\end{equation}
mit $O$,$R$ den Austrittswellen von Objekt und Referenz. Für die aufgenommene Intensität $I$ im Fernfeld gilt
\begin{equation}
I=\left|\mathscr{F}\left[O\right]+\mathscr{F}\left[R\right]\right|^2=\left|\tilde{O}+\tilde{R}\right|^2
\end{equation}
Wird nun auf $I$ eine inverse Fouriertransformation angewendet,

\begin{equation}
	\mathscr{F}^{-1}\left[\left|\tilde{O}+\tilde{R}\right|^2\right]=
	\mathscr{F}^{-1}\left[\tilde{O}\tilde{O}^*\right]+
	\mathscr{F}^{-1}\left[\tilde{R}\tilde{R}^*\right]+
	\mathscr{F}^{-1}\left[\tilde{R}\tilde{O}^*\right]+
		\mathscr{F}^{-1}\left[\tilde{O}\tilde{R}^*\right]
\end{equation}
so ergibt sich mit Hilfe der komplexen Konjugation der Fouriertransformation XX und des Korrelationstheorems XX
\begin{equation}
\mathscr{F}^{-1}[I]\propto \underbrace{O \otimes O + R\otimes R}_{\text{Autokorrelationen}}+\underbrace{R\otimes O+ O\otimes R}_{\text {Kreuzkorrelationen}}
\end{equation}
Die Autokorrelationen, d.h. die Korrelation von Objekt und Objekt bzw Referenz mit Referenz liegen hierbei um den Ursprung des Koordinatensystems und können in jede Richtung maximal die doppelte Größe von Objekt bzw Referenz haben. Die Kreuzkorrelationen aus Objekt und Referenz sind punktsymmetrisch zum Ursprung und haben von diesem den gleichen Abstand wie Objekt- und Referenzaustrittswelle voneinander.
Sie  entsprechen einer Faltung des Objektes mit der gespiegelten, komplex konjugierten Referenz bzw der Referenz mit einem gespiegelten, komplex konjugierten Objekt.  Ist die Referenz in guter Näherung punktförmig und lässt sich durch eine Delta-Funktion beschreiben, so ist die Kreuzkorrelation näherungsweise die Austrittswelle des Objektes. Somit können trotz der verlorenen Phaseninformationen die Austrittswelle des Objektes inklusiver der Phase wieder gewonnen werden. Dieses Verfahren wird als im weiteren als Fouriertransformations-Holografie (FTH) bezeichnet.
\subsection{Entfaltung}
Um bei ausgedehnten Referenzen eine bessere Approximation des Objektes zu erhalten, muss die Faltung zwischen Objekt und Referenz rückgängig gemacht werden. Nach  XXX entspricht die Faltung einer Multiplikation, die Inverse somit einer Division im Fourierraum. Jedoch verursachen sowohl Rauschen wie auch Singulariäten in der fouriertransformierten Referenz Probleme.

Das Problem der Entfaltung tritt in verschiedenen Bereichen auf, unter anderem in der Bildverarbeitung.
Hier ist der Prozess der Wiener Entfaltung etabliert. 

Für Rauschen, dass nicht mit dem Signal korreliert ist stellt dies die Lösung...

\section{iterative Phasenrückgewinnung}
Wäre neben der Amplitude auch Phase der Fouriertransformation der Austrittswelle bekannt, so könnte diese durch eine inverse Transformation wiederhergestellt werden. 
Da bei N komplexen Datenpunkten im Realraum die Fouriertransformation durch N Amplituden und N Phasen, also 2N Variablen dargestellt werden kann, jedoch nur N Variablen (die Fourieramplituden) gemessen werden können, ist das Problem unterbestimmt und nicht eindeutig lösbar. Kann jedoch das zu rekonstruierende Objekt im Realraum auf N/2 Bildpunkte beschränkt werden, so wird das Problem lösbar. Die räumliche Beschränkung im Realraum wird 
durch die Wahl eines sog. Supports, außerhalb dessen der Realraum als leer angenommen wird umgesetzt.  
Die Problemstellung ist nun eine Lösung der Gleichung $\rho(x)=\mathscr{F}^{-1}\tilde{\rho}(q)$ unter den Nebenbedingungen dass 1. $\rho$ auf den Support beschränkt ist und 2. die Amplituden von $\tilde{\rho}$ mit den bekannten Fourieramplituden $A(q)$ übereinstimmt. 
\subsection{Algorithmen}
Ein zur Lösung der Problemstellung verwendeter Algorithmus, der Error-Reduction Algorithmus (ER) basiert auf abwechselnden Erzwingen der der Nebenbedingungen, d.h. von einem Startwert $\rho_0$ ausgehend wird zunächst $\tilde{\rho_0}$ bestimmt, die Amplitude im Fourierraum durch $A$ ersetzt, in den Realraum transformiert und dort außerhalb des Supports gleich Null gesetzt. Dies wird iterativ wiederholt, wobei sich $\rho_n$ einer beide Nebenbedingungen erfüllenden Lösung annähert.

Das Erzwingen der Nebenbedinung im Realraum kann als Projektion auf die Lösungsmenge $S$ der Nebenbedingung  mit dem Komplement $\bar{S}$ durch den Projektor $P_s$ 

 \begin{align}
 P_s\rho (x)=\begin{cases}
 \rho (x)  &\text{für } x\in S\\
 \mathbb{0}  &\text{für }x\notin S
 \end{cases}&&
 P_{\bar{s}}\rho (x)=\begin{cases}
 \mathbb{0} &\text{für } x\in S\\
 \rho (x)   &\text{für }x\notin S
 \end{cases}
 \end{align}
aufgefasst werden. Der analoge Projektor für die Erfüllung der Nebenbedingung im Fourierraum, der im Fourierraum die Amplitude von $\tilde{\rho}$ durch die bekannte Amplitude $A(q)$ ersetzt und dabei den Winkel $\Phi(q)$ beibehält lautet
\begin{align}
	\tilde{P}_m \bar{\rho}(q)&=A(q)e^{i\Phi(q)} &&P_m\rho(q)=\mathscr{F}^{-1}A(q)e^{i\Phi(q)}\mathscr{F}\rho(x)
\end{align}

	\begin{figure}
		\centering
		\includegraphics[width=0.75\textwidth]{images/projektor.eps}
		\caption[Abbildungstext im Abbildungsverzeichnis]{Verwendete Projektoren. Der Support Projektor $P_S$ projiziert auf die durch den Support beschriebene Mannigfaltigkeit, $P_{\bar{S}}$ auf das Komplement. Ein Reflektor wendet die doppelte Schrittweite wie ein Projektor an und stellt somit eine Spiegelung an der Lösungsmenge der Nebenbedingung dar. Der Fourier-Projektor stellt in Fourierraum eine projektion auf die Kreisförmige Menge der bekannten Amplituden dar.}
		\label{fig:projektoren}
	\end{figure}.
	Neben den Projektoren die, den Schritt $\rho\rightarrow (\mathbb{1}+(P-\mathbb{1}))\rho$ ausführen, können auch sogenannte Reflektoren mit doppelter Schrittweite
	\begin{equation}
	R_\nu= (\mathbb{1}+2(P-\mathbb{1}))=(2P_\nu-\mathbb{1})
	\end{equation}
	eingeführt werden.
In dieser Schreibweise können auch Weiterentwicklungen wie der Hybrid Input Output (HIO) und Relaxed Averaged Alternating Reflections (RAAR) Algorithmus dargestellt werden XXTabelleXX. Ihre Prinzip ist in XXX graphisch dargestellt.
	\begin{figure}
		\centering
		\includegraphics[width=0.75\textwidth]{images/algorithmen.eps}
		\caption[Abbildungstext im Abbildungsverzeichnis]{Darstellung der Algorithmen ER, RAAR und HIO. Bei ER findet eine abwechselnde Projektion auf die Lösungsmengen der Nebenbedinungen statt. Bei ER und HIO ist der einzelne Iterationsschritt komplexer und abhängig von einem Parameter $\beta$. Abbildung für ER und HIO nach \cite{XX}}
		\label{Abb:BspAbbildung}
	\end{figure} 
\subsection{Support}
Üblicherweise ist der Bereich auf den das Objekt im Realraum beschränkt ist nicht bekannt, der für die iterativen Rekonstruktionsalgorithmen benötigten Supports muss zunächst bestimmt werden. Hierfür existieren verschiedene Möglichkeiten. Die triviale Methode besteht aus einer Festlegung aus zuvor bekannten Dimensionen des Experiments. Jedoch kann nur in Ausnahmefällen so der Support ausreichend eng gewählt werden um die Bedingung XXX zur erfüllen.
Eine verbreitete Methode, bezeichnet als Shrink-Wrap Algorithmus, ist durch eine Anpassung des Supports während der Rekonstruktion gekennzeichnet: Der zunächst geratene (für eine eindeutige Lösung des Problems  zu große) Support wird nach einer festgelegten Anzahl an Rekonstruktionsiterationen auf den Bereich beschränkt, in dem die Rekonstruierte Intensität höher als ein festgelegter Schwellwert ist. Dieser aktualisierte Support wird für weitere Iterationen genutzt und erneut aktualisiert\cite{marchesini2003}.
Die aus der Holographie gewonnene Rekonstruktion des Objektes und eine Schätzung der Referenz lassen sich zu einem Support kombinieren. Der nötige Abstandsvektor dieser Supportteile ist der  Abstandsvektor zwischen Auto- und Kreuzkorrelation (XXX).

Abbildung
	\begin{figure}
		\centering
		\includegraphics[width=0.75\textwidth]{images/support.eps}
		\caption[Abbildungstext im Abbildungsverzeichnis]{Erzeugung des Supports. Oben: Bei verwendung von Freiflug Holographie kann die Kreuzkorrelation aus Referenz und Objekt als Support genutzt werden. Zusätzlich muss im gleichen Abstand wie Kreuz und Autokorrelation die Referenz dem Support hinzugefügt werden. Unten: Der Shrinkwrap Alorithmus beginnt mit der Autokorrelation als initalem Support und passt diesen im laufe der Rekonstruktion an.}
		\label{Abb:BspAbbildung}
	\end{figure} 
\section{Vergleich}	
Somit können drei unterschiedliche Arten der Rekonstruktion betrachtet werden
\begin{itemize}
\item FTH mit Entfaltung
\item IPR mit Shrinkwrap-Support
\item IPR mit FTH Support
\end{itemize}

Die Rekonstruktionsansätze werden bezüglich Ihrer Empfindlichkeit für Rauschen, Hochpass und Fehler in der Abschätzung der Referenz untersucht. Dies geschieht zunächst mit rein synthetischen 2D Bildern als Ersatz für Austrittswellen um die Funktion der drei Ansätze und ihre Limitationen zu überprüfen.

\subsection{Einfluss Rauschen}
Zur Beurteilung des Einflusses von Rauschen auf die Rekonstruktionen werden die Fourieramplituden mit verschiedenen Auflösungen diskretisiert und ein Poisson-Rauschen angewendet.
\subsection{Einfluss Hochpass}
Der Einfluss des in Experimenten durch die nötige Blockierung des zentralen Strahles bedingten Hochpasses
\subsection{Einfluss Referenzabschätzung}
Sowohl für die Entfaltung wie auch für die festlegung des Supports für die iterative Rekonstruktion muss die Referenz abgeschätzt werden. Der Einfluss der Güte dieser Abschätzung..

\section{3D Austrittswelle}
Die für 2D-Bilder getesteten Ansätze können nun zur Rekonstruktion der Austrittswelle hinter einem 3D-Objekt eingesetzt werden. Die Dreidimensionalität des Objektes wird an dieser Stelle ignoriert und die Austrittswelle analog zu den 2D-Bildern betrachtet.
\section{Ausblick}
